\section{Method}
\label{sec:method}

\subsection{Problem Formulation}
\label{sec:formulation}

We formulate per-head KV cache quantization as a constrained optimization problem.
Let $h \in \{1,\ldots,H\}$ index the KV heads across all layers, $b_h \in \mathcal{B}$ the bit-width assigned to head $h$, and $\mathcal{B} = \{2, 4, 8, 16\}$ the set of hardware-aligned bit-widths.
Each head $h$ has a profiled distortion function $d_h(b)$, the PPL increase when head $h$ is quantized to $b$ bits.
Given a memory budget ratio $R$ (e.g., $R = 0.25$ for $4\times$ compression), we seek the allocation that minimizes total distortion:
\begin{equation}
\label{eq:objective}
\min_{\{b_h\}} \sum_{h=1}^{H} d_h(b_h) \quad \text{s.t.} \quad \frac{1}{H \cdot 16}\sum_{h=1}^{H} b_h \leq R.
\end{equation}

This is a multiple-choice knapsack problem (MCKP): each head selects one bit-width from $\mathcal{B}$ to minimize total distortion under a budget.
RateQuant~\cite{zuo2026ratequant} solves this for consecutive bit-widths $\{2,3,4,5,6\}$, where the polymatroid structure of the feasible set guarantees that greedy allocation is optimal (their Proposition 6).
Hardware-aligned bit-widths $\{2,4,8,16\}$ are non-consecutive; the gaps break this structure, and greedy is no longer guaranteed to match the MCKP optimum.
Moreover, as we show in \S\ref{sec:compounding}, the objective itself is not well-behaved: distortions from 2-bit heads compound non-additively, making unconstrained optimization unreliable.

\subsection{PPL-based Per-Head Profiling}
\label{sec:profiling}

We use symmetric per-channel quantization with per-head granularity: each KV head is independently quantized to one of the candidate bit-widths $\mathcal{B} = \{2, 4, 8, 16\}$, without per-token scaling.
We profile each head by quantizing it to each $b \in \mathcal{B}$ while keeping all other heads at bf16, and measuring the resulting PPL increase $d_h(b)$ on a calibration set of 64 sequences $\times$ 2048 tokens from WikiText-2 train.
This yields a distortion table of $H \times |\mathcal{B}|$ entries that serves as input to the solver.

The optimization in Eq.~\eqref{eq:objective} relies on an additive distortion assumption: the total PPL degradation from a mixed-precision allocation is approximately $\sum_h d_h(b_h)$.
This holds when per-head distortions are small and independent---a reasonable approximation for 4-bit and 8-bit quantization, where individual head distortions are typically below 0.1 PPL.
At 2-bit, however, individual distortions are orders of magnitude larger, and the additive assumption breaks down.

\subsection{Compounding Effect}
\label{sec:compounding}

When many heads are simultaneously quantized to 2-bit, per-head distortions compound non-additively: truncation errors from multiple 2-bit heads accumulate through the residual stream, and once their count exceeds the network's capacity to compensate, PPL degrades super-linearly.

Unconstrained solvers are blind to this effect because the additive objective underestimates the cost of large-scale 2-bit assignment.
On Qwen2.5-7B, unconstrained greedy assigns 68 heads to 2-bit and achieves 32.45 PPL; unconstrained MCKP assigns 75 heads and achieves 36.92---both far worse than uniform 4-bit (15.45).
Standard-resolution profiling detects non-zero $d_h(2)$ for only 37 of 112 heads; the remaining 75 heads have $d_h(2) \approx 0$, so most 2-bit assignments appear free to the additive objective.
The resulting additive prediction is $\sum_h d_h(b_h) + \text{baseline} \approx 11$ PPL---$3\text{--}4\times$ lower than the actual 32--37 PPL, because the compounding interaction between simultaneously quantized heads is invisible to per-head profiling.

This compounding is specific to 2-bit quantization and is \emph{super-additive}: quantizing the top-$k$ most sensitive heads to 2-bit simultaneously produces PPL $1.56\times$ the sum of their individual distortions at $k{=}8$, growing to $4.89\times$ at $k{=}64$.
A random-$k$ control confirms that sensitivity ordering matters: at $k{=}64$, random selection yields 55.94 PPL vs.\ 22.95 for top-$k$, because random draws are more likely to include high-distortion heads that amplify compounding.
At 4-bit and above, per-head distortions are small enough that additivity holds even when all heads share the same bit-width (uniform 4-bit achieves 15.45 PPL, close to the sum-of-distortions prediction).
The qualitative difference arises because 2-bit quantization has far larger per-head distortion, saturating the error tolerance of the residual stream.

The relationship between $k$ and PPL is U-shaped (Figure~\ref{fig:compounding}, Appendix~\ref{sec:compounding_appendix}): at $k{=}0$ (uniform 4-bit), PPL is 15.45; allowing a few 2-bit heads reduces PPL by freeing budget for 8-bit heads on sensitive positions, reaching a minimum of 10.91 at $k{=}8$; beyond $k{=}8$, compounding drives PPL upward.

\subsection{Cardinality Constraint}
\label{sec:cardinality}

The compounding effect motivates a simple addition to Eq.~\eqref{eq:objective}: a cardinality constraint on 2-bit assignments.
\begin{equation}
\label{eq:constraint}
\sum_{h=1}^{H} \mathbb{1}[b_h = 2] \leq k.
\end{equation}

The constraint parameter $k$ is selected by a small-scale sweep on calibration data: solve the constrained problem for each candidate $k$, evaluate calibration PPL, and pick the value at the minimum of the U-shaped curve (Figure~\ref{fig:compounding}).
For Qwen2.5-7B with 112 KV heads, $k{=}8$ ($\sim$7\% of heads).

The constraint is a single linear inequality, compatible with any solver---greedy, MCKP, or integer linear programming---and adds negligible overhead to the optimization.

One might expect Layer-0 heads to require explicit protection.
On Qwen2.5-7B, however, the profiler assigns very high 2-bit distortion to Layer-0 heads (L0H0: $+$31.9\%, L0H2: $+$23.2\%, L0H3: $+$11.7\% relative to bf16), so no solver would assign them 2-bit under any reasonable budget.
Running the solver with and without manual Layer-0 exclusion produces identical allocations; no manual rule is needed.

\paragraph{Restoring additivity.}
The additive objective treats per-head distortions as independent, but 2-bit compounding violates this once many 2-bit heads coexist.
Restricting 2-bit heads to $k \ll H$ pushes the allocation back into a near-additive regime where the profiler's prediction agrees with measured PPL.
For Qwen2.5-7B, $k{=}8$ ($\sim$7\% of heads) is few enough that compounding contributes negligibly, consistent with the U-shaped curve minimum (Figure~\ref{fig:compounding}).
This explains the constrained--unconstrained gap in Table~\ref{tab:solver_equiv}: constrained greedy operates inside the additive regime (10.91 PPL, matching the ${\approx}11$ additive estimate), while unconstrained solvers extrapolate outside it and miss true PPL by $3\text{--}4\times$.

\subsection{Solver Equivalence}
\label{sec:solver_equiv}

\begin{table}[t]
\centering
\small
\caption{Solver comparison on Qwen2.5-7B at $4\times$ compression. Unconstrained solvers over-allocate 2-bit heads and fail catastrophically. The cardinality constraint ($k{\leq}8$) makes greedy and MCKP produce identical allocations.\protect\footnotemark}
\label{tab:solver_equiv}
\begin{tabular}{llcccc}
\toprule
Solver & Constraint & PPL & \multicolumn{3}{c}{Bit distribution} \\
 & & & 2-bit & 4-bit & 8/16-bit \\
\midrule
Greedy & None & 32.45 & 68 & 30 & 14 \\
MCKP & None & 36.92 & 75 & 33 & 4 \\
Greedy & $k{\leq}8$ & 10.91 & 8 & 100 & 4 \\
MCKP & $k{\leq}8$ & 10.91 & 8 & 100 & 4 \\
Uniform & --- & 15.45 & 0 & 112 & 0 \\
\bottomrule
\end{tabular}
\end{table}
\footnotetext{Constrained results use high-resolution profiling; unconstrained rows use standard-resolution profiling to illustrate the failure mode (the exact PPL is irrelevant---any unconstrained allocation fails catastrophically).}

With the cardinality constraint, greedy and MCKP produce identical results (Table~\ref{tab:solver_equiv}): same PPL (10.91), same bit distribution on Qwen2.5-7B.

\paragraph{Empirical observation.}
When $k \ll H$ and the top-$k$ heads ranked by marginal benefit $d_h(4) - d_h(2)$ contain no ties, constrained greedy and MCKP produce identical per-head allocations.
The constraint restricts the feasible set so that the only decision is which $k$ heads to assign 2-bit; the marginal-benefit ordering determines this uniquely, so greedy's local choice coincides with the global optimum.
We verify this across three architectures: Qwen2.5-7B (112/112 heads identical), Llama-3.1-8B (256/256), Mistral-7B-v0.3 (255/256; the single difference is a zero-distortion tie with no PPL impact).
Without the constraint, greedy cascades into excessive 2-bit allocation, as each additional 2-bit head looks locally beneficial but contributes to compounding.
Practitioners can therefore use the simpler $O(H \log H)$ greedy algorithm without sacrificing optimality.

A per-layer brute-force oracle on Qwen2.5-7B confirms near-optimality: of 8 layers with non-trivial allocations, MCKP matches the oracle on 5; the remaining 3 have a maximum gap of 0.018 PPL (0.4\% of the total improvement), closing the optimality chain from greedy through MCKP to brute-force.